\chapter{Evaluation of Results}
\label{ch:evaluation}

This section evaluates the performance of the classification and clustering components with a focus on interpreting the results rather than re-describing the methodology. The goal is to understand how the observed metrics reflect the structure of the dataset and the behaviour of the \gls{emma} pipeline as a whole. 

\section{Classification Performance and Interpretation}
The final classification model achieves a weighted F1-score of 0.69 and a Cohen's $\kappa$ of 0.66 on a held-out validation set. These values are computed after retraining the selected \gls{tfidf} Bigrams + \gls{lsvc} model on the full MedMCQA corpus and evaluating on a 10\% holdout split. The champion configuration, \gls{tfidf} Bigrams + \gls{lsvc}, was identified through 10-fold cross-validation (Section \ref{sec:classification_model_selection}) and retrained on the full 179,777-question corpus for holdout evaluation. 

During model selection, this configuration achieves a weighted F1-score of 0.5424 $\pm$ 0.0086 and a Cohen's $\kappa$ of 0.5089 $\pm$ 0.0096, obtained through 10-fold cross-validation. The final performance increase to 0.69 / 0.66 is explained by retraining the selected configuration on the full dataset, which significantly increases the amount of available training data. 

A key factor influencing classification performance is the structure of the dataset itself. This is quantified through inter-category similarity analysis, where \gls{tfidf} centroid cosine similarity between specialty classes is computed. The observed mean similarity of approximately 0.7101 indicates that categories are moderately distinct from one another. This level of separability makes it feasible for a linear classifier to learn effective decision boundaries, while still presenting a non-trivial multi-class problem due to the presence of 19 specialties.

The dominance of \gls{tfidf} bigrams over embedding-based representations can be directly understood from this dataset structure. Specialty labels in MedMCQA are strongly associated with domain-specific terminology, meaning that short lexical patterns (e.g., ``myocardial infarction'', ``gestational diabetes'') act as highly discriminative signals. \gls{tfidf} bigrams explicitly capture these local phrase-level patterns, whereas dense embeddings tend to smooth over such distinctions by emphasizing global semantic similarity. As a result, embedding-based models underperform in this task, as reflected in the grid search outputs. 

\section{Clustering Performance and Interpretation}
The key insight from the clustering results is that topic-level structure in medical questions is more granular than specialty-level categorization. While classification assigns each question to one of 19 specialties, clustering reveals subtopics within those specialties, such as cardiac symptoms within Internal Medicine or gestational conditions within Obstetrics. This difference in granularity explains the low $\kappa$ scores and highlights the complementary roles of classification and clustering in \gls{emma}. Classification provides coarse-grained routing, while clustering provides fine-grained refinement. In the retrieval pipeline, this allows the system to first narrow the search space to a relevant specialty and then further restrict it to a topic-specific subset of documents. 

The outlier rate of 34.7\% is also consistent with the characteristics of the dataset. MedQA questions are relatively short (approximately 20 words on average), which limits the density structure required by \gls{hdbscan} to form stable clusters. Rather than discarding these outliers, the system falls back to specialty-based routing for these cases, ensuring that no queries are lost. From a systems perspective, clustering enhances retrieval precision by introducing an additional layer of structure between classification and retrieval. Even when cluster assignments are imperfect, they still guide the retrieval process toward more relevant regions of the corpus, improving the overall quality of the \gls{rag} pipeline. 

\section{NER as a Retrieval Quality Gate}
The \gls{ner} analysis (Chapter \ref{ch:ner_collocation}) reveals that retrieval quality is embedding-dependent. \gls{ner} rewriting improves Octen-Embedding scores by +0.006 but hurts MiniLM by $-0.022$. This explains run 2 vs. run 6 in the ablation grid: the same \gls{rag} condition produces a $-4$\gls{pp} result with one embedding and a +11\gls{pp} result with another. The \gls{ner} component is necessary but not sufficient as its benefit is only realised when the downstream embedding can interpret biomedical entity strings. 

\section{System-Level Implications}
Classification and clustering serve complementary roles that collectively improve retrieval precision beyond what either achieves alone. Classification provides a stable coarse-grained specialty label that narrows the retrieval search space; clustering then introduces finer-grained topic grouping within that specialty. Questions assigned to the outlier cluster (topic = $-1$) retain their specialty label and fall back to specialty-only routing without any loss of functionality. The \gls{ner} component (Chapter \ref{ch:ner_collocation}) adds a third layer: by rewriting the retrieval query to entity strings before \gls{faiss} search, it addresses embedding dilution that classification and clustering cannot. Together these three deterministic, auditable components constitute the routing layer of \gls{emma}, leaving the \gls{llm} to focus exclusively on explanation generation.
